\subsection{Completion of the critical transition}
The accumulated remainder conditions in Theorem~\ref{thm:v7-critical}
are useful for heterogeneous finite schedules. For a homogeneous
sequence a condition on the error of one observation suffices throughout
the entire transition, even when the full accumulated error diverges.

\begin{theorem}[Zero, finite, and infinite effective sample limits]
\label{thm:v8-supercritical}
Let $j\to\infty$, $d_j>0$, $q_j=e^{-j\gamma}$, and let $\omega$ have
the meaning in Lemma~\ref{lem:v7-tangent}. The geometry may vary with
$j$ in a fixed compact family; $\gamma,A$ and the laws then take
their corresponding parameter values. Suppose
\begin{equation}\label{eq:v8-one-error}
                     \omega(d_j)/q_j\longrightarrow0.
\end{equation}
If $k_jq_j\to b\in[0,\infty]$, then
\[
 \|\Pi_{j,d_j}^{\otimes k_j}-
       \Pi_{\partial,j,d_j}^{\otimes k_j}\|_{\rm TV}
                               \longrightarrow1-e^{-2b/\pi}.
\]
If $n_jp_{j,d_j}^0q_j\to b\in[0,\infty]$, the same limit holds for
$\|P_{j,d_j}^{\otimes n_j}-P_{\partial,j,d_j}^{\otimes n_j}\|_{\rm TV}$.
At $b=\infty$ the right side is one. The minimum equal-prior errors
converge to $e^{-2b/\pi}/2$. Both statements permit unequal facing
curvatures and either parity. Under \eqref{eq:v8-one-error}, vanishing
total variation is equivalent to $k_jq_j\to0$ in the conditional
experiment, and to $n_jp_{j,d_j}^0q_j\to0$ in the raw experiment.
\end{theorem}
\begin{proof}
If $b<\infty$, then
$k_j\omega(d_j)=(k_jq_j)\omega(d_j)/q_j\to0$ in the conditional
case, and
$n_jp_{j,d_j}^0\omega(d_j)=(n_jp_{j,d_j}^0q_j)\omega(d_j)/q_j\to0$
in the raw case. Theorem~\ref{thm:v7-critical} applies, including $b=0$.

Suppose $b=\infty$ and fix $B>0$. For all sufficiently large $j$ the
conditional sample contains its first $\lfloor B/q_j\rfloor$
observations. Projection onto them cannot increase total variation.
Their effective sample size tends to $B$ and their accumulated tangent
error tends to zero by \eqref{eq:v8-one-error}. Hence the lower limit
of the full-sample variation distance is at least $1-e^{-2B/\pi}$.
Let $B\to\infty$ and use the upper bound one. For raw preparations
project instead onto the first $\lfloor B/(p_{j,d_j}^0q_j)\rfloor$
observations. Here $p_{j,d_j}^0q_j\to0$, so rounding changes their
effective sample size by a quantity tending to zero. The identical
argument proves the raw assertion. All estimates used above are uniform on the compact geometric family,
and its positive lower bound for $\gamma$ ensures $q_j\to0$ even
when the geometry varies. If an effective sample size fails to tend to
zero, a subsequence is bounded below by a positive constant. Project
on that subsequence onto a fixed positive critical subsample smaller
than this constant. Its limiting distance is positive by the finite
case, proving necessity in the final equivalences. Sufficiency is the
case $b=0$. The testing formula follows from
$\inf(\text{equal-prior error})=(1-\|P-Q\|_{\rm TV})/2$.
\end{proof}

For a general smooth contact, \eqref{eq:v8-one-error} is
$d_j=o(q_j^2)$; when both graphs are even it is $d_j=o(q_j)$.
The finite critical subsamples therefore approach their quadratic
experiments in total variation. This is a sharp physical testing
profile in an asymptotically tangent regime, not a replacement for the
nonlinear fixed-offset theorem or a proof that its rate $\tau$ is sharp.
The supercritical conclusion no longer requires the tangent error of
all retained observations to vanish. The subsample argument applies to
the known simple pair of laws, and does not supply a simulator for an
unknown table.
